\section{Introduction and Problem Understanding}
Air pollution from fine particulate matter with aerodynamic diameter under 2.5 micrometers ($\text{PM}_{2.5}$) poses major respiratory and cardiovascular health hazards. Exposure to elevated $\text{PM}_{2.5}$ levels increases the incidence of chronic bronchitis, asthma, and premature mortality. In South Africa, physical air quality monitoring stations are sparse and unevenly distributed, leaving several municipalities without continuous ground observations. Accurate data-driven forecasting of particulate matter enables municipal authorities to issue timely health advisories and plan mitigation policies.

In this Phase 1 project, we implemented and evaluated the baseline machine learning methodology published by Morapedi and Obagbuwa (2023) \cite{morapedi2023air}, who applied gradient boosting algorithms and standard classifiers to historical air quality records across twelve South African cities. Our objective was to build a complete, independent Python pipeline that reproduces their data processing, continuous regression modeling, spatial cross-validation, temporal evaluation, future forecasting, and Air Quality Index (AQI) classification. In the following sections, we summarize the original paper, describe the dataset, detail our preprocessing pipeline, outline our models and experimental setup, present our empirical findings, compare our numbers directly with the published paper, and discuss methodological flaws identified in the original study. We record that the DOI link provided in the course assignment sheet (\texttt{https://doi.org/10.3390/s21041064}) points to a different paper (MDPI Sensors 21(4):1064), whereas the paper title assigned corresponds to the Frontiers in Artificial Intelligence article by Morapedi and Obagbuwa \cite{morapedi2023air}, which we implemented.
